\documentclass[11pt]{amsart}

\usepackage[margin=1in]{geometry}
\usepackage{amsmath,amssymb,amsthm,mathtools}
\usepackage{enumitem}

\newtheorem{theorem}{Theorem}
\newtheorem{proposition}[theorem]{Proposition}
\newtheorem{lemma}[theorem]{Lemma}
\newtheorem{corollary}[theorem]{Corollary}
\newtheorem{remark}[theorem]{Remark}
\newtheorem{definition}[theorem]{Definition}

\newcommand{\E}{\mathbb{E}}
\newcommand{\C}{\mathbb{C}}
\newcommand{\N}{\mathbb{N}}
\newcommand{\R}{\mathbb{R}}
\newcommand{\T}{\mathbb{T}}
\newcommand{\supp}{\operatorname{supp}}
\newcommand{\gapgc}{\operatorname{gap}_{\mathrm{gc}}}
\newcommand{\dT}{d_{\mathbb{T}}}

\title{Complex stable phase retrieval with one exceptional coordinate}

\author{Pedro Abdalla Teixeira\textsuperscript{1}}
\author{Jaume de Dios Pont\textsuperscript{2}}
\author{Jo\~ao Pedro Ramos\textsuperscript{3}}
\author{Mitchell Taylor\textsuperscript{4}}

\begin{document}
\begin{abstract}
We give a compactness proof of the complex one-exception stable phase retrieval
theorem.  The proof is organized in the same head-tail architecture as the real
one-exception argument: one first proves the coefficient decompositions and the
head identification, then extracts a limit model in which a head pair and a
tail pair are independent and satisfy an exact modulus identity, and finally
proves a support theorem for the tail difference.  Two complementary support
routes are recorded for the degenerate head case.  The first is purely
geometric: cone-supported infinitely divisible laws are supported on additive
subgroups of the cone, hence on either a complex line or a conjugate plane, and
both possibilities force the tail difference onto a real line.  The second,
included as a remark after the proof, uses the nontriviality of the Gaussian
part of a diffuse tail limit.  Along the way we isolate the quantitative
angular-spread estimate needed for the complex cross-coordinate probes and we
import the rigorously proved compactness ingredients from the general complex
stable phase retrieval theorem.
\end{abstract}
\maketitle
\begingroup
\renewcommand{\thefootnote}{\arabic{footnote}}
\footnotetext[1]{Department of Mathematics, University of California, Irvine,
  Irvine, California 92697, USA}
\footnotetext[2]{Center for Data Science, New York University, New York,
  New York 10011, USA}
\footnotetext[3]{Instituto de Matem\'atica Pura e Aplicada (IMPA), Rio de
  Janeiro, RJ, Brazil}
\footnotetext[4]{Department of Mathematics, ETH Z\"urich, Z\"urich, Switzerland}
\endgroup
\setcounter{footnote}{0}

\section{Introduction}

Let $(\Omega,\mu)$ be a probability space and let $\xi_1,\ldots,\xi_n$ be
independent complex-valued random variables.  The main theorem of
\cite{complex_spr} states that if every $\xi_j$ satisfies a two-sided
$L^1$-$L^2$ norm-equivalence bound and a uniform lower bound on the
generalized-circle gap $\gapgc(\xi_j)\ge \kappa>0$, then there is a uniform
orthogonal modulus gap
\[
  \bigl\|\,|X_a|-|X_b|\,\bigr\|_{L^2} \ge c(A,B,\kappa)>0
\]
for every pair of orthogonal unit vectors $a,b\in\C^n$.

The theorem proved here is the one-exception analogue of
\cite[Theorem~2.1]{complex_spr}: a single exceptional coordinate $\eta$ is
allowed to violate the lower $L^1$-$L^2$ bound and the generalized-circle gap,
while the good coordinates retain the same uniform hypotheses as in the fully
homogeneous theorem.

The present proof keeps the current \texttt{.tex} file as the main source but
merges into it the rigorously proved compactness ingredients from
Appendix~A of \cite{complex_spr_app}.  The resulting organization mirrors the
real one-exception proof from \cite{compactness_real}.  First one proves the
head-tail decompositions and the head identification by bounded probes.  Next
one packages those lemmas into a compactness reduction, whose output is an
exact modulus identity
\[
  |H_X+R_X|=|H_Y+R_Y|
\]
for an independent head pair and tail pair.  The proof then culminates in a
support theorem for the tail difference.  When the head is nontrivial, the
support theorem follows from a geometric elimination of the head variable.  In
the degenerate case, one uses the additive-subgroup structure of cone-supported
infinitely divisible laws to force the tail difference onto a real line.

For completeness, we also keep the Gaussian route to the degenerate case.  That
route is not needed in the final proof, but it gives an independent way to see
that the same support conclusion follows from the nontriviality of the Gaussian
part of a diffuse tail limit.

\section{Background and notation}
\label{sec:background}

\begin{definition}[Generalized circle]
A \emph{generalized circle} in $\C$ is either a Euclidean circle
$\{z \in \C : |z-c|=r\}$ with $c \in \C$ and $r \ge 0$, or a real affine line
in $\C \cong \R^2$.
\end{definition}

\begin{definition}[Generalized-circle gap {\cite[Definitions~3.3--3.4]{complex_spr}}]
\label{def:gapgc}
For a complex-valued random variable $\xi \in L^2(\mu)$, define
\[
  \gapgc(\xi)
  :=
  \inf_{\substack{A,C\in\R,\;B\in\C \\ A^2+2|B|^2+C^2=1}}
  \bigl\|A|\xi|^2 + 2\,\Re(B\xi) + C\bigr\|_{L^1}.
\]
By \cite[Lemma~3.8]{complex_spr}, $\gapgc(\xi)=0$ if and only if
$\supp(\xi)$ is contained in a generalized circle.
\end{definition}

\begin{definition}[Phase-invariant distance {\cite[Definition~2.2]{complex_spr}}]
For $f,g \in L^2(\mu)$, set
\[
  \dT(f,g) := \inf_{\lambda \in \T} \|f-\lambda g\|_{L^2(\mu)}.
\]
\end{definition}

We use the following lemmas from \cite{complex_spr}.

\begin{lemma}[One-coordinate rigidity {\cite[Lemma~3.10]{complex_spr}}]
\label{lem:one-coord}
Let $S \subset \C$ be a set not contained in any generalized circle.  Let
$\alpha,\beta,u,v \in \C$ and assume
\[
  |\alpha t + u| = |\beta t + v| \qquad \text{for every } t \in S.
\]
Then $|\alpha|=|\beta|$ and $\bar\alpha\,u=\bar\beta\,v$.  Consequently, either
$\alpha=\beta=0$, or there exists $\lambda \in \T$ such that
$\beta=\lambda\alpha$ and $v=\lambda u$.
\end{lemma}

\begin{lemma}[Support of an independent pair {\cite[Lemma~3.11]{complex_spr}}]
\label{lem:supp-prod}
If $U$ and $V$ are independent random variables taking values in Polish spaces,
then $\supp(U,V)=\supp(U)\times \supp(V)$.
\end{lemma}

\begin{lemma}[Almost-sure vanishing extends to the support
  {\cite[Lemma~3.12]{complex_spr}}]
\label{lem:as-support}
Let $S$ be a Polish space, let $Z$ be an $S$-valued random variable, and let
$F:S\to\R$ be continuous.  If $F(Z)=0$ almost surely, then
$F(x)=0$ for every $x\in\supp(Z)$.
\end{lemma}

\section{Main result and contradiction setup}
\label{sec:main}

\begin{theorem}[Complex one-exception orthogonal lower bound]
\label{thm:main}
Let $\eta$ be a complex-valued random variable satisfying
\[
  \E[\eta]=0,
  \qquad
  \E[|\eta|^2]=1.
\]
Fix constants $0 < A \le B < 1$ and $\kappa > 0$.  Then there exists
\[
  c = c(A,B,\kappa,\eta) > 0
\]
such that the following holds for every $n \in \N$.

Let $\xi_1,\ldots,\xi_n$ be complex-valued random variables, independent of
each other and of $\eta$, satisfying
\[
  \E[\xi_j]=0,
  \qquad
  \E[|\xi_j|^2]=1
  \qquad (1\le j\le n),
\]
and assume:
\begin{enumerate}[label=\rm(\arabic*)]
\item for every $u=(u_1,\ldots,u_n)\in\C^n$,
\[
  A\|u\|_2
  \le
  \Bigl\|\sum_{j=1}^n u_j\xi_j\Bigr\|_{L^1}
  \le
  B\|u\|_2;
\]
\item $\gapgc(\xi_j)\ge \kappa$ for every $1\le j\le n$.
\end{enumerate}
For every $a=(a_0,a_1,\ldots,a_n)$ and $b=(b_0,b_1,\ldots,b_n)$ in $\C^{n+1}$
with
\[
  \|a\|_2=\|b\|_2=1,
  \qquad
  \langle a,b\rangle=0,
\]
define
\[
  X_a := a_0\eta + \sum_{j=1}^n a_j\xi_j,
  \qquad
  X_b := b_0\eta + \sum_{j=1}^n b_j\xi_j.
\]
Then
\[
  \bigl\|\,|X_a|-|X_b|\,\bigr\|_{L^2} \ge c.
\]
\end{theorem}

\begin{remark}
Theorem~\ref{thm:main} is the natural one-exception analogue of
\cite[Theorem~2.1]{complex_spr}.  The constant is allowed to depend on the full
law of the exceptional variable $\eta$.
\end{remark}

\begin{corollary}[Complex one-exception stable phase retrieval]
\label{cor:spr}
Under the hypotheses of Theorem~\ref{thm:main}, let
$E_n:=\operatorname{span}\{\eta,\xi_1,\ldots,\xi_n\}\subset L^2(\mu)$.  There
exists $C=C(A,B,\kappa,\eta)<\infty$ such that for all $f,g\in E_n$,
\[
  \dT(f,g) \le C\,\bigl\|\,|f|-|g|\,\bigr\|_{L^2(\mu)}.
\]
\end{corollary}

\begin{proof}
This is proved exactly as in \cite[Corollary~2.4]{complex_spr} once the
orthogonal lower bound is known.
\end{proof}

\begin{proposition}[Normalized counterexample array]
\label{prop:counterexample}
If Theorem~\ref{thm:main} fails, then there exist:
\begin{enumerate}[label=\rm(\roman*)]
\item a triangular array $(\xi_{n,i})_{n\ge 1,\ i\ge 1}$ such that for each
  fixed $n$ the family
  \[
    (\eta,\xi_{n,1},\xi_{n,2},\dots)
  \]
  is independent, every $\xi_{n,i}$ is centered and satisfies
  \[
    \E[|\xi_{n,i}|^2]=1,
    \qquad
    \gapgc(\xi_{n,i})\ge \kappa,
  \]
  and for every finitely supported $u\in\ell^2(\N)$,
  \[
    A\|u\|_2
    \le
    \Bigl\|\sum_{i\ge 1} u_i\xi_{n,i}\Bigr\|_{L^1}
    \le
    B\|u\|_2;
  \]
\item vectors $a^{(n)},b^{(n)}\in \ell^2(\{0,1,2,\dots\})$ such that
  \[
    \|a^{(n)}\|_2=\|b^{(n)}\|_2=1,
    \qquad
    \langle a^{(n)},b^{(n)}\rangle=0,
  \]
  and, writing
  \[
    X_n:=a_0^{(n)}\eta+\sum_{i\ge 1} a_i^{(n)}\xi_{n,i},
    \qquad
    Y_n:=b_0^{(n)}\eta+\sum_{i\ge 1} b_i^{(n)}\xi_{n,i},
  \]
  one has
  \[
    \bigl\|\,|X_n|-|Y_n|\,\bigr\|_{L^2}\longrightarrow 0.
  \]
\end{enumerate}
\end{proposition}

\begin{proof}
This is the standard contradiction setup: one takes violating orthogonal unit
pairs row by row, and pads finite vectors by zeros.
\end{proof}

\section{Coefficient decompositions and the limit model}
\label{sec:compactness}

We now formulate the compactness scheme in the same language as
\cite{compactness_real}.  The point is that the proof should work directly with
the bad sequence of sums rather than first extracting limit laws for the
individual coordinates.

\begin{lemma}[Reordered profile decomposition]
\label{lem:profile}
After passing to a subsequence in Proposition~\ref{prop:counterexample} and, in
each row, applying the same permutation to the good coordinates and to their
coefficients, one may find vectors
\[
  a,b\in \ell^2(\{0,1,2,\dots\})
\]
and integers $m_n\to\infty$ such that, with
\[
  a^{(n)}_{\mathrm{head}}:=a^{(n)}\mathbf 1_{\{0,1,\dots,m_n\}},
  \qquad
  a^{(n)}_{\mathrm{tail}}:=a^{(n)}\mathbf 1_{\{m_n+1,m_n+2,\dots\}},
\]
and likewise for $b^{(n)}$, the following hold:
\begin{enumerate}[label=\rm(\arabic*)]
\item $a^{(n)}_{\mathrm{head}}\to a$ and
  $b^{(n)}_{\mathrm{head}}\to b$ strongly in $\ell^2$;
\item
\[
  \max_{i>m_n}\max\bigl(|a_i^{(n)}|,|b_i^{(n)}|\bigr)\longrightarrow 0;
\]
\item
\[
  \|a^{(n)}_{\mathrm{tail}}\|_2^2\longrightarrow 1-\|a\|_2^2,
  \qquad
  \|b^{(n)}_{\mathrm{tail}}\|_2^2\longrightarrow 1-\|b\|_2^2,
\]
and
\[
  \langle a^{(n)}_{\mathrm{tail}},b^{(n)}_{\mathrm{tail}}\rangle
  \longrightarrow
  -\langle a,b\rangle.
\]
\end{enumerate}
\end{lemma}

\begin{proof}
The proof is the same monotone-reordering and diagonal extraction used in the
real one-exception paper; it depends only on the coefficient vectors.
\end{proof}

\begin{remark}
The exceptional coordinate $\eta$ is kept fixed at index $0$.  Only the good
coordinates are reordered.
\end{remark}

\begin{proposition}[Bounded probes for good coordinates]
\label{prop:bounded-probes}
Let $\xi$ be a centered complex-valued random variable satisfying
\[
  \E[|\xi|^2]=1,
  \qquad
  A\le \|\xi\|_{L^1}\le B<1.
\]
Then there exist bounded measurable functions $S,T:\Omega\to\C$ such that
\[
  \E[S]=\E[T]=0,
  \qquad
  \E[\xi\overline{S}]=\E[\overline{\xi}S]=1,
  \qquad
  \E[(|\xi|^2-1)T]=1,
\]
and
\[
  \|S\|_{L^\infty}\le \frac{2}{A},
  \qquad
  \|T\|_{L^\infty}\le \frac{2}{1-B}.
\]
\end{proposition}

\begin{proof}
Define the phase function
\[
  \phi :=
  \begin{cases}
    \xi/|\xi|,& \xi\ne 0,\\
    0,& \xi=0.
  \end{cases}
\]
Then $|\phi|\le 1$ and
\[
  \E[\xi\overline{\phi}] = \E[|\xi|] = \|\xi\|_{L^1}\ge A.
\]
Set
\[
  S:=\frac{\phi-\E[\phi]}{\E[|\xi|]}.
\]
Since $\E[\xi]=0$, one has
\[
  \E[\xi\overline{S}]
  =
  \frac{\E[\xi\overline{\phi}]-\E[\xi]\E[\overline{\phi}]}{\E[|\xi|]}
  = 1.
\]
Taking complex conjugates gives $\E[\overline{\xi}S]=1$, and clearly
$\E[S]=0$.  The bound
\[
  \|S\|_{L^\infty}\le \frac{2}{\E[|\xi|]}\le \frac{2}{A}
\]
is immediate.

Next let
\[
  \psi:=\operatorname{sign}(|\xi|^2-1),
\]
which is real-valued and bounded by $1$ in modulus.  Then
\[
  \E[(|\xi|^2-1)\psi]=\E\bigl[\,||\xi|^2-1|\,\bigr].
\]
Moreover,
\[
  ||\xi|^2-1| = ||\xi|-1|(|\xi|+1)\ge ||\xi|-1|,
\]
so
\[
  \E\bigl[\,||\xi|^2-1|\,\bigr]
  \ge
  \E[\,||\xi|-1|\,]
  \ge
  1-\E[|\xi|]
  \ge
  1-B.
\]
Therefore, with
\[
  T:=\frac{\psi-\E[\psi]}{\E[(|\xi|^2-1)\psi]},
\]
we get $\E[T]=0$,
\[
  \E[(|\xi|^2-1)T]=1,
\]
and
\[
  \|T\|_{L^\infty}
  \le
  \frac{2}{\E[\,||\xi|^2-1|\,]}
  \le
  \frac{2}{1-B}.
\]
\end{proof}

\begin{lemma}[Quantitative angular spread]
\label{lem:angular-spread}
Let $\xi$ be a centered complex-valued random variable with
\[
  \E[|\xi|^2]=1.
\]
Define
\[
  m:=\E[|\xi|],
  \qquad
  c:=\E\!\left[\frac{\xi^2}{|\xi|}\right],
\]
where $\xi^2/|\xi|$ is interpreted as $0$ on $\{\xi=0\}$.  Then
\[
  |c|\le m-\gapgc(\xi)^2.
\]
In particular, if $S$ is the linear probe from
Proposition~\ref{prop:bounded-probes}, then
\[
  \rho:=\E[\xi S]=\frac{c}{m}
\]
satisfies
\[
  |\rho|\le 1-\gapgc(\xi)^2\le 1.
\]
\end{lemma}

\begin{proof}
If $c=0$, there is nothing to prove.  Otherwise choose $\omega\in\T$ so that
\[
  \omega^2 c = |c|.
\]
Write $\omega\xi = |\xi|e^{it}$.  Then
\[
  |c|
  =
  \Re(\omega^2 c)
  =
  \E\!\left[|\xi|\cos(2t)\right],
\]
so
\[
  m-|c|
  =
  \E\!\left[|\xi|\bigl(1-\cos(2t)\bigr)\right]
  =
  2\,\E\!\left[|\xi|\sin^2 t\right].
\]

Consider the generalized circle given by the line
\[
  \Im(\omega z)=0.
\]
The corresponding normalized test has the form
\[
  z\longmapsto \sqrt2\,\Im(\omega z),
\]
and therefore
\[
  \gapgc(\xi)\le \sqrt2\,\E\!\left[|\Im(\omega\xi)|\right]
  =
  \sqrt2\,\E\!\left[|\xi||\sin t|\right].
\]
By Cauchy--Schwarz,
\[
  \E\!\left[|\xi||\sin t|\right]^2
  \le
  \E[|\xi|]\,
  \E\!\left[|\xi|\sin^2 t\right]
  =
  \frac{m(m-|c|)}{2}.
\]
Combining the last two displays gives
\[
  \gapgc(\xi)^2\le m(m-|c|)\le m-|c|,
\]
because $m\le \|\xi\|_{L^2}=1$.  This proves the first claim.  The identity
\[
  \E[\xi S]
  =
  \frac{\E[\xi^2/|\xi|]}{\E[|\xi|]}
  =
  \frac{c}{m}
\]
follows directly from the definition of $S$, and the stated bound on $|\rho|$
is immediate.
\end{proof}

\begin{proposition}[Linear probe for the exceptional coordinate]
\label{prop:exceptional-probe}
Let $\eta$ be a centered complex-valued random variable satisfying
\[
  \E[|\eta|^2]=1.
\]
Then there exists a bounded measurable function $S_\eta$ such that
\[
  \E[S_\eta]=0,
  \qquad
  \E[\eta\overline{S_\eta}]
  =
  \E[\overline{\eta}S_\eta]
  = 1.
\]
\end{proposition}

\begin{proof}
Define
\[
  \phi_\eta :=
  \begin{cases}
    \eta/|\eta|,& \eta\ne 0,\\
    0,& \eta=0.
  \end{cases}
\]
Then $\E[|\eta|]>0$ because $\E[|\eta|^2]=1$.  Set
\[
  S_\eta:=\frac{\phi_\eta-\E[\phi_\eta]}{\E[|\eta|]}.
\]
Exactly as above, $\E[S_\eta]=0$ and
\[
  \E[\eta\overline{S_\eta}]
  =
  \E[\overline{\eta}S_\eta]
  =1.
\]
\end{proof}

\begin{proposition}[Head identification]
\label{prop:head}
Assume the setup of Lemma~\ref{lem:profile}.  Then either:
\begin{enumerate}[label=\rm(\arabic*)]
\item there exists $\tau\in\T$ such that
\[
  a=\tau b
\]
and
\[
  \|a^{(n)}_{\mathrm{head}}-\tau b^{(n)}_{\mathrm{head}}\|_2\longrightarrow 0;
\]
\item or else
\[
  a_i=b_i=0
  \qquad\text{for every } i\ge 1.
\]
\end{enumerate}
\end{proposition}

\begin{proof}
Since
\[
  \bigl\||X_n|^2-|Y_n|^2\bigr\|_{L^1}
  \le
  \bigl\||X_n|-|Y_n|\bigr\|_{L^2}
  \bigl\||X_n|+|Y_n|\bigr\|_{L^2}
  \le
  2\bigl\||X_n|-|Y_n|\bigr\|_{L^2},
\]
we have
\[
  \bigl\||X_n|^2-|Y_n|^2\bigr\|_{L^1}\longrightarrow 0.
\]

Fix $i\ge 1$.  Let $T_{n,i}$ be the quadratic probe given by
Proposition~\ref{prop:bounded-probes} for $\xi_{n,i}$.  Since $T_{n,i}$ is a
bounded centered function of $\xi_{n,i}$ alone, it is independent of $\eta$ and
of all $\xi_{n,j}$ with $j\ne i$.  Hence
\[
  \E\bigl[(|X_n|^2-|Y_n|^2)T_{n,i}\bigr]\longrightarrow 0.
\]
Expanding $|X_n|^2=X_n\overline{X_n}$ and using independence and centering, all
terms vanish except the one involving $(|\xi_{n,i}|^2-1)T_{n,i}$.  Therefore
\[
  \E\bigl[(|X_n|^2-|Y_n|^2)T_{n,i}\bigr]
  =
  |a_i^{(n)}|^2-|b_i^{(n)}|^2.
\]
Passing to the limit gives
\[
  |a_i|^2=|b_i|^2
  \qquad (i\ge 1).
\]

Now fix distinct indices $i,j\ge 1$.  Let $S_{n,i}$ and $S_{n,j}$ be the linear
probes from Proposition~\ref{prop:bounded-probes}.  Since
$\overline{S_{n,i}}S_{n,j}$ is bounded and depends only on
$(\xi_{n,i},\xi_{n,j})$, we again have
\[
  \E\bigl[(|X_n|^2-|Y_n|^2)\overline{S_{n,i}}S_{n,j}\bigr]\longrightarrow 0.
\]
Set
\[
  A_{n,i,j}:=
  a_i^{(n)}\overline{a_j^{(n)}}-b_i^{(n)}\overline{b_j^{(n)}},
  \qquad
  \rho_{n,\ell}:=\E[\xi_{n,\ell}S_{n,\ell}]
  \qquad (\ell=i,j).
\]
As before, independence and centering eliminate every term except the two
containing $\xi_{n,i}\overline{\xi_{n,j}}$ and $\xi_{n,j}\overline{\xi_{n,i}}$.
Using
\[
  \E[\xi_{n,i}\overline{S_{n,i}}]=1,
  \qquad
  \E[\overline{\xi_{n,j}}S_{n,j}]=1,
\]
and taking complex conjugates where needed, we obtain
\[
  \E\bigl[(|X_n|^2-|Y_n|^2)\overline{S_{n,i}}S_{n,j}\bigr]
  =
  A_{n,i,j}
  +
  \overline{A_{n,i,j}}\,\rho_{n,j}\,\overline{\rho_{n,i}}.
\]
By Lemma~\ref{lem:angular-spread} and the standing lower bound
$\gapgc(\xi_{n,\ell})\ge\kappa$, one has
\[
  |\rho_{n,\ell}|
  \le
  1-\kappa^2
  \qquad (\ell=i,j).
\]
Hence
\[
  \bigl|A_{n,i,j}\bigr|
  \le
  \frac{
    \bigl|\E\bigl[(|X_n|^2-|Y_n|^2)\overline{S_{n,i}}S_{n,j}\bigr]\bigr|
  }{
    1-(1-\kappa^2)^2
  }
  \longrightarrow 0,
\]
because
\[
  \bigl|A_{n,i,j}
  +
  \overline{A_{n,i,j}}\,\rho_{n,j}\,\overline{\rho_{n,i}}\bigr|
  \ge
  \bigl(1-|\rho_{n,j}\rho_{n,i}|\bigr)\,|A_{n,i,j}|.
\]
Passing to the limit yields
\[
  a_i\overline{a_j}=b_i\overline{b_j}
  \qquad (i\ne j,\ i,j\ge 1).
\]

Let $S_\eta$ be given by Proposition~\ref{prop:exceptional-probe}.  For every
$j\ge 1$, letting $S_{n,j}$ be the good-coordinate linear probe for $\xi_{n,j}$,
we likewise obtain
\[
  \E\bigl[(|X_n|^2-|Y_n|^2)\overline{S_\eta}S_{n,j}\bigr]\longrightarrow 0
\]
and therefore
\[
  a_0^{(n)}\overline{a_j^{(n)}}-b_0^{(n)}\overline{b_j^{(n)}}\longrightarrow 0.
\]
Passing to the limit gives
\[
  a_0\overline{a_j}=b_0\overline{b_j}
  \qquad (j\ge 1).
\]

Assume first that there exists $i_0\ge 1$ with $a_{i_0}\ne 0$.  Since
$|a_{i_0}|=|b_{i_0}|$, there exists $\tau\in\T$ such that
\[
  a_{i_0}=\tau b_{i_0}.
\]
For every $j\ge 1$,
\[
  a_{i_0}\overline{a_j}
  =
  b_{i_0}\overline{b_j}
  =
  \overline{\tau}a_{i_0}\overline{b_j},
\]
so $a_j=\tau b_j$.  Taking $j=i_0$ in the exceptional-good identity gives
\[
  a_0\overline{a_{i_0}}
  =
  b_0\overline{b_{i_0}}
  =
  \tau b_0\overline{a_{i_0}},
\]
and therefore $a_0=\tau b_0$.  Thus $a=\tau b$.

If instead $a_i=0$ for every $i\ge 1$, then $|a_i|=|b_i|$ implies
$b_i=0$ for every $i\ge 1$, which is exactly alternative \textup{(2)}.

In alternative \textup{(1)}, strong convergence of the head vectors gives
\[
  \|a^{(n)}_{\mathrm{head}}-\tau b^{(n)}_{\mathrm{head}}\|_2
  \le
  \|a^{(n)}_{\mathrm{head}}-a\|_2
  +
  \|b^{(n)}_{\mathrm{head}}-b\|_2
  \longrightarrow 0.
\]
\end{proof}

\begin{remark}
The linear probe in the complex case is exactly the phase analogue of the sign
probe in the real case: one replaces $\operatorname{sign}(\xi)$ by
$\xi/|\xi|$.  The quadratic probe is unchanged, since $|\xi|^2-1$ is real.
\end{remark}

\begin{lemma}[Limit law for the head-tail decomposition]
\label{lem:limit-law}
With the notation above, set
\[
  X_n^{\mathrm{head}}
  :=
  a_0^{(n)}\eta+\sum_{1\le i\le m_n} a_i^{(n)}\xi_{n,i},
  \qquad
  X_n^{\mathrm{tail}}
  :=
  \sum_{i>m_n} a_i^{(n)}\xi_{n,i},
\]
and define $Y_n^{\mathrm{head}},Y_n^{\mathrm{tail}}$ analogously.  After
passing to a further subsequence, the quadruple
\[
  \bigl(
    X_n^{\mathrm{head}},
    Y_n^{\mathrm{head}},
    X_n^{\mathrm{tail}},
    Y_n^{\mathrm{tail}}
  \bigr)
\]
converges in distribution to
\[
  (H_X,H_Y,R_X,R_Y)\in\C^4,
\]
where:
\begin{enumerate}[label=\rm(\arabic*)]
\item $(H_X,H_Y)$ is independent of $(R_X,R_Y)$;
\item
\[
  |H_X+R_X|=|H_Y+R_Y|
  \qquad\text{almost surely;}
\]
\item either there exists $\tau\in\T$ such that
  \[
    H_X=\tau H_Y
    \qquad\text{almost surely,}
  \]
  or else
  \[
    H_X=a_0\eta,
    \qquad
    H_Y=b_0\eta
    \qquad\text{almost surely.}
  \]
\end{enumerate}
\end{lemma}

\begin{proof}
The tightness and independence statements are the same as in the real proof.
The last alternative is just Proposition~\ref{prop:head} transported to random
variables: in alternative \textup{(1)} of Proposition~\ref{prop:head}, strong
$L^2$ convergence of the head coefficient vectors gives
\[
  X_n^{\mathrm{head}}-\tau Y_n^{\mathrm{head}}\longrightarrow 0
\]
in $L^2$, hence $H_X=\tau H_Y$ almost surely.  In alternative \textup{(2)},
\[
  \sum_{1\le i\le m_n} |a_i^{(n)}|^2
  \longrightarrow 0,
  \qquad
  \sum_{1\le i\le m_n} |b_i^{(n)}|^2
  \longrightarrow 0,
\]
so
\[
  X_n^{\mathrm{head}}-a_0\eta\longrightarrow 0,
  \qquad
  Y_n^{\mathrm{head}}-b_0\eta\longrightarrow 0
\]
in $L^2$.
\end{proof}

\begin{proposition}[Compactness reduction to an exact head-tail identity]
\label{prop:compactness-reduction}
Assume that Theorem~\ref{thm:main} fails, and let the bad sequence be given by
Proposition~\ref{prop:counterexample}.  Then, after passing to a subsequence,
there exist:
\begin{enumerate}[label=\rm(\arabic*)]
\item an integer sequence $m_n\to\infty$ and reordered coefficient vectors
  $a^{(n)},b^{(n)}$ as in Lemma~\ref{lem:profile};
\item complex random variables $(H_X,H_Y,R_X,R_Y)$ such that
  $(H_X,H_Y)$ is independent of $(R_X,R_Y)$;
\item limiting head coefficients $a=(a_0,a_1,\dots)$ and $b=(b_0,b_1,\dots)$
  satisfying the alternative of Proposition~\ref{prop:head};
\item the exact modulus identity
  \[
    |H_X+R_X|=|H_Y+R_Y|
    \qquad\text{almost surely.}
  \]
\end{enumerate}
Moreover, either there exists $\tau\in\T$ such that
\[
  H_X=\tau H_Y
  \qquad\text{almost surely,}
\]
or else
\[
  H_X=a_0\eta,
  \qquad
  H_Y=b_0\eta
  \qquad\text{almost surely.}
\]
\end{proposition}

\begin{proof}
This is exactly the content of Lemma~\ref{lem:profile},
Proposition~\ref{prop:head}, and Lemma~\ref{lem:limit-law}.
\end{proof}

\begin{lemma}[Lower $L^1$-$L^2$ comparability for independent real sums]
\label{lem:real-l1l2}
Let $\zeta_1,\dots,\zeta_N$ be independent centered real-valued random
variables and assume that
\[
  \|\zeta_i\|_{L^1}\ge \delta \|\zeta_i\|_{L^2}
  \qquad (1\le i\le N)
\]
for some $\delta>0$.  Then
\[
  \Bigl\|\sum_{i=1}^N \zeta_i\Bigr\|_{L^1}
  \ge
  \frac{\delta}{2\sqrt 2}
  \Bigl(\sum_{i=1}^N \|\zeta_i\|_{L^2}^2\Bigr)^{1/2}.
\]
\end{lemma}

\begin{proof}
If $\sum_i \|\zeta_i\|_{L^2}^2=0$, there is nothing to prove.  By homogeneity, we
may therefore assume
\[
  \sum_{i=1}^N \|\zeta_i\|_{L^2}^2=1.
\]
Let $\zeta_i'$ be an independent copy of $\zeta_i$, and set
\[
  S:=\sum_{i=1}^N \zeta_i,
  \qquad
  S':=\sum_{i=1}^N \zeta_i'.
\]
Then
\[
  2\|S\|_{L^1}\ge \|S-S'\|_{L^1}
  =
  \Bigl\|\sum_{i=1}^N (\zeta_i-\zeta_i')\Bigr\|_{L^1}.
\]
The differences $\zeta_i-\zeta_i'$ are independent and symmetric.  Conditioning
on their absolute values and applying the Khintchine inequality gives
\[
  \Bigl\|\sum_{i=1}^N (\zeta_i-\zeta_i')\Bigr\|_{L^1}
  \ge
  \frac{1}{\sqrt 2}
  \E\Bigl[\Bigl(\sum_{i=1}^N (\zeta_i-\zeta_i')^2\Bigr)^{1/2}\Bigr].
\]
Since
\[
  \sum_{i=1}^N c_i^2 |\widetilde\zeta_i-\widetilde\zeta_i'|
  \le
  \Bigl(\sum_{i=1}^N c_i^2(\widetilde\zeta_i-\widetilde\zeta_i')^2\Bigr)^{1/2},
\]
where we write
\[
  c_i:=\|\zeta_i\|_{L^2},
  \qquad
  \widetilde\zeta_i :=
  \begin{cases}
    \zeta_i/c_i,& c_i\ne 0,\\
    0,& c_i=0,
  \end{cases}
\]
so that $\sum_i c_i^2=1$ and $\|\widetilde\zeta_i\|_{L^2}=1$ whenever $c_i\ne 0$.
Applying the previous argument to the sum $\sum_i c_i \widetilde\zeta_i$ gives
\[
  \|S\|_{L^1}
  \ge
  \frac{1}{2\sqrt 2}
  \sum_{i=1}^N c_i^2
  \|\widetilde\zeta_i-\widetilde\zeta_i'\|_{L^1}.
\]
Finally,
\[
  \|\widetilde\zeta_i-\widetilde\zeta_i'\|_{L^1}
  \ge
  \Bigl\|\E[\widetilde\zeta_i-\widetilde\zeta_i'\mid \widetilde\zeta_i]\Bigr\|_{L^1}
  =
  \|\widetilde\zeta_i\|_{L^1}
  \ge \delta,
\]
because $\widetilde\zeta_i'$ is centered and
$\|\widetilde\zeta_i\|_{L^1}\ge \delta \|\widetilde\zeta_i\|_{L^2}=\delta$.
Hence
\[
  \Bigl\|\sum_{i=1}^N \zeta_i\Bigr\|_{L^1}
  \ge
  \frac{\delta}{2\sqrt 2}.
\]
Undoing the normalization yields the stated inequality.
\end{proof}

\begin{lemma}[Infinitely divisible laws supported on a generalized circle are line-supported]
\label{lem:id-gc-line}
Let $\mu$ be an infinitely divisible probability measure on $\C$.  If
$\supp\mu$ is contained in a generalized circle, then $\supp\mu$ is contained in
an affine line.
\end{lemma}

\begin{proof}
If $\supp\mu$ is already contained in an affine line, there is nothing to
prove.  It remains to exclude the case in which $\supp\mu$ is contained in a
Euclidean circle $\Gamma$.

Because $\mu$ is infinitely divisible, it is in particular $2$-divisible:
there exists a probability measure $\nu$ such that $\mu=\nu*\nu$.  Let
$A:=\supp\nu$.  Then
\[
  \supp\mu = \overline{A+A}.
\]
Assume for contradiction that $A$ contains two distinct points $x\ne y$.  Then
\[
  2x,\ x+y,\ 2y \in A+A \subset \supp\mu \subset \Gamma.
\]
But $x+y$ is the midpoint of the chord joining $2x$ and $2y$.  The midpoint of a
nontrivial chord of a Euclidean circle lies strictly inside the enclosed disk,
and therefore cannot belong to the circle itself.  This contradiction shows
that $A$ is a singleton.  Hence $\mu=\nu*\nu$ is a point mass, and in
particular its support is contained in an affine line.
\end{proof}

\begin{proposition}[Generalized-circle tail rigidity]
\label{prop:tail-rigidity}
Let $(\xi_{n,i})_{n\ge 1,\ i\ge 1}$ be a triangular array of independent
centered complex-valued random variables satisfying
\[
  \E[|\xi_{n,i}|^2]=1,
  \qquad
  \gapgc(\xi_{n,i})\ge \kappa,
\]
and the uniform finite-support $L^1$-$L^2$ bounds
\[
  A\|u\|_2
  \le
  \Bigl\|\sum_{i\ge 1} u_i\xi_{n,i}\Bigr\|_{L^1}
  \le
  B\|u\|_2.
\]
Let $c^{(n)}\in\ell^2(\N)$ satisfy
\[
  \sup_n\|c^{(n)}\|_2<\infty,
  \qquad
  \max_i |c_i^{(n)}|\longrightarrow 0,
\]
and define
\[
  Z_n:=\sum_{i\ge 1} c_i^{(n)}\xi_{n,i}.
\]
If $Z_n$ converges in distribution to a random variable $Z$ whose law is
supported on a generalized circle, then
\[
  \|c^{(n)}\|_2\longrightarrow 0.
\]
\end{proposition}

\begin{proof}
Assume for contradiction that $\|c^{(n)}\|_2\not\to 0$.  After passing to a
subsequence, we may assume that
\[
  \delta \le \|c^{(n)}\|_2 \le M
  \qquad\text{for every } n
\]
for some constants $\delta>0$ and $M<\infty$.

Since
\[
  \E[|Z_n|^2]=\sum_{i\ge 1} |c_i^{(n)}|^2 = \|c^{(n)}\|_2^2 \le M^2,
\]
the family $(Z_n)_n$ is bounded in $L^2$.  In particular, the first moments are
uniformly integrable.  Because every $Z_n$ is centered, passing to the weak
limit gives
\[
  \E[Z]=0.
\]

The array of summands
\[
  Z_{n,i}:=c_i^{(n)}\xi_{n,i}
\]
is infinitesimal: for every $\varepsilon>0$,
\[
  \sup_i \mathbb P\bigl(|Z_{n,i}|>\varepsilon\bigr)
  \le
  \frac{1}{\varepsilon^2}\sup_i |c_i^{(n)}|^2
  \longrightarrow 0.
\]
By the classical theorem on weak limits of row-wise independent infinitesimal
triangular arrays, the law of $Z$ is infinitely divisible.  Since by hypothesis
it is also supported on a generalized circle, Lemma~\ref{lem:id-gc-line} shows
that $\supp(Z)$ is contained in an affine line.  Because $\E[Z]=0$, that line
passes through the origin.  Hence there exists $\lambda\in\T$ such that
\[
  \lambda Z \in \R
  \qquad\text{almost surely.}
\]

Define
\[
  Y_n := \Im(\lambda Z_n)
  = \sum_{i\ge 1} Y_{n,i},
  \qquad
  Y_{n,i}:=\Im(\lambda c_i^{(n)}\xi_{n,i}).
\]
Then each $Y_{n,i}$ is a centered real random variable, the family
$(Y_{n,i})_i$ is independent for each fixed $n$, and
\[
  Y_n \Rightarrow 0
\]
in distribution.

We claim that each nonzero summand $Y_{n,i}$ satisfies a uniform lower
$L^1$-$L^2$ comparability estimate.  If $c_i^{(n)}\ne 0$, write
\[
  \omega_{n,i} := \lambda\,\frac{c_i^{(n)}}{|c_i^{(n)}|}\in\T.
\]
Then
\[
  Y_{n,i} = |c_i^{(n)}|\,\Im(\omega_{n,i}\xi_{n,i}).
\]
For every $\omega\in\T$, choose
\[
  A=0,\qquad C=0,\qquad B=-\,\frac{i\overline{\omega}}{\sqrt 2}.
\]
Then
\[
  A^2+2|B|^2+C^2=1
\]
and
\[
  A|\xi|^2+2\Re(B\xi)+C = \sqrt 2\,\Im(\omega\xi).
\]
Therefore
\[
  \kappa
  \le
  \gapgc(\xi_{n,i})
  \le
  \sqrt 2\,\|\Im(\omega_{n,i}\xi_{n,i})\|_{L^1}
  \le
  \sqrt 2\,\|\Im(\omega_{n,i}\xi_{n,i})\|_{L^2}.
\]
Multiplying by $|c_i^{(n)}|$ gives
\[
  \|Y_{n,i}\|_{L^1} \ge \frac{\kappa}{\sqrt 2}|c_i^{(n)}|,
  \qquad
  \|Y_{n,i}\|_{L^2} \ge \frac{\kappa}{\sqrt 2}|c_i^{(n)}|.
\]
Since also $\|Y_{n,i}\|_{L^2}\le |c_i^{(n)}|$, we get
\[
  \|Y_{n,i}\|_{L^1}
  \ge
  \frac{\kappa}{\sqrt 2}\|Y_{n,i}\|_{L^2}.
\]

Applying Lemma~\ref{lem:real-l1l2} row by row, we obtain
\[
  \|Y_n\|_{L^1}
  \ge
  \frac{\kappa}{4}
  \Bigl(\sum_{i\ge 1}\|Y_{n,i}\|_{L^2}^2\Bigr)^{1/2}
  \ge
  \frac{\kappa^2}{4\sqrt 2}\|c^{(n)}\|_2
  \ge
  \frac{\kappa^2\delta}{4\sqrt 2}.
\]

On the other hand, $Y_n\Rightarrow 0$ and
\[
  \sup_n \E[|Y_n|^2] \le \sup_n \E[|Z_n|^2]\le M^2,
\]
so the family $(|Y_n|)_n$ is uniformly integrable.  Hence
\[
  \|Y_n\|_{L^1}\longrightarrow 0,
\]
which contradicts the uniform positive lower bound above.
\end{proof}

\begin{remark}
Proposition~\ref{prop:tail-rigidity} is the complex replacement for the
one-dimensional tail rigidity used in the real proof.  Its proof separates the
generalized-circle obstruction into two parts: infinite divisibility rules out
the curved case, and the remaining line case reduces to the real
$L^1$-$L^2$ argument after taking imaginary parts in the correct direction.
\end{remark}

\section{Tail support lemmas}
\label{sec:proof}

\begin{lemma}[Geometric reduction when the head is nontrivial]
\label{lem:nontrivial-head-geometry}
Let $H$ be a complex-valued random variable, let $(R_X,R_Y)$ be a
complex-valued random vector independent of $H$, and let $\tau\in\T$.  Assume
that
\[
  |H+R_X| = |\tau H + R_Y|
  \qquad\text{almost surely.}
\]
If $\supp(H)$ contains at least two distinct points, then the law of
\[
  W:=R_Y-\tau R_X
\]
is supported on an affine line.
\end{lemma}

\begin{proof}
By independence and Lemmas~\ref{lem:supp-prod} and~\ref{lem:as-support}, the
identity
\[
  |H+R_X| = |\tau H + R_Y|
\]
holds for every
\[
  h\in \supp(H),
  \qquad
  (x,y)\in \supp(R_X,R_Y).
\]
Choose two distinct points $h_1,h_2\in \supp(H)$.  Fix
$(x,y)\in \supp(R_X,R_Y)$ and write
\[
  w:=y-\tau x.
\]
For $j=1,2$ one has
\[
  |h_j+x| = |\tau(h_j+x)+w|.
\]
Squaring and cancelling $|h_j+x|^2$ gives
\[
  |w|^2 + 2\Re\bigl(\overline{\tau(h_j+x)}\,w\bigr)=0
  \qquad (j=1,2).
\]
Subtracting the two identities yields
\[
  \Re\bigl(\overline{\tau(h_1-h_2)}\,w\bigr)=0.
\]
Thus every point of the image of $\supp(R_X,R_Y)$ under the continuous map
\[
  (x,y)\mapsto y-\tau x
\]
lies in the affine line
\[
  L:=\{w\in\C:\ \Re(\overline{\tau(h_1-h_2)}\,w)=0\}.
\]
Since $L$ is closed, the support of the pushforward law of $W$ is also
contained in $L$.
\end{proof}

\begin{remark}
This section collects the support lemmas used in the final compactness
argument.  The nontrivial-head case is handled by a direct geometric
elimination of the head variable.  The degenerate case $H_X\equiv H_Y\equiv 0$
is treated by classifying additive subgroups of the cone
\[
  \Gamma:=\{(u,v)\in\C^2:\ |u|=|v|\}.
\]
\end{remark}

\begin{lemma}[Two distinct unimodular slopes force real-line degeneracy]
\label{lem:two-slopes-real-degenerate}
Let $\tau_1,\tau_2\in\T$ with $\tau_1\ne \tau_2$, and for $j=1,2$ let
\[
  L_{\tau_j}:=\{(z,\tau_j z): z\in\C\}\subset \C^2.
\]
Let $S\subset (L_{\tau_1}\cup L_{\tau_2})\cap \{(u,v)\in\C^2:\ |u|=|v|\}$ and
assume that
\[
  S+S\subset \{(u,v)\in\C^2:\ |u|=|v|\}.
\]
For $j=1,2$, define
\[
  K_j:=\{z\in\C:\ (z,\tau_j z)\in S\}.
\]
Then the following hold.
\begin{enumerate}[label=\rm(\arabic*)]
\item For every $z_1\in K_1$ and $z_2\in K_2$,
  \[
    \Re\bigl((1-\tau_1\overline{\tau_2})z_1\overline{z_2}\bigr)=0.
  \]
\item If $K_2$ contains a nonzero point, then $K_1$ is contained in a real line
  through the origin.
\item If $K_1$ and $K_2$ both contain nonzero points, then both sets are
  contained in real lines through the origin.
\end{enumerate}
\end{lemma}

\begin{proof}
Fix $z_1\in K_1$ and $z_2\in K_2$.  Then
\[
  x:=(z_1,\tau_1 z_1)\in S,
  \qquad
  y:=(z_2,\tau_2 z_2)\in S.
\]
By hypothesis, $x+y\in S+S$ belongs to the cone, so
\[
  |z_1+z_2| = |\tau_1 z_1+\tau_2 z_2|.
\]
Squaring and using $|\tau_1|=|\tau_2|=1$, we obtain
\[
  |z_1|^2+|z_2|^2+2\Re(z_1\overline{z_2})
  =
  |z_1|^2+|z_2|^2+2\Re(\tau_1\overline{\tau_2}\,z_1\overline{z_2}),
\]
hence
\[
  \Re\bigl((1-\tau_1\overline{\tau_2})z_1\overline{z_2}\bigr)=0.
\]
This proves \textup{(1)}.

Assume now that $z_2\in K_2$ is nonzero.  Since $\tau_1\ne \tau_2$, the complex
number
\[
  \lambda:=(1-\tau_1\overline{\tau_2})\overline{z_2}
\]
is nonzero.  By \textup{(1)}, every $z_1\in K_1$ satisfies
\[
  \Re(\lambda z_1)=0,
\]
which describes a real line through the origin in $\C\cong \R^2$.  This proves
\textup{(2)}.  Interchanging the roles of the two slopes gives \textup{(3)}.
\end{proof}

\begin{lemma}[Covariance range of a non-line-supported rank-one complex vector]
\label{lem:cov-range-rank-one}
Let $\xi$ be a centered complex-valued random variable with
\[
  \E[|\xi|^2]<\infty,
\]
and let $(a,b)\in\C^2\setminus\{(0,0)\}$.  Define
\[
  V:=(a\xi,b\xi)\in\C^2.
\]
Assume that $\xi$ is not supported on a real line in $\C$.  Then the range of
the real covariance operator of $V$ is exactly the real plane underlying the
complex line
\[
  L:=\{(az,bz): z\in\C\}.
\]
\end{lemma}

\begin{proof}
Since $V$ is supported on $L$, the covariance range is contained in the
underlying real plane $L_\R$.

Conversely, let $\ell:\C^2\to\R$ be a real linear functional in the kernel of
the covariance operator.  Because $V$ is centered, this means
\[
  \E[\ell(V)^2]=0,
\]
hence $\ell(V)=0$ almost surely.  By continuity, $\ell$ vanishes on
$\supp(V)$.  Since $\xi$ is not supported on a real line, the real span of
$\supp(\xi)$ is all of $\C$, and therefore the real span of $\supp(V)$ is the
full plane $L_\R$.  Thus $\ell$ vanishes on $L_\R$, so the annihilator of the
covariance range equals the annihilator of $L_\R$.  Hence the covariance range
is exactly $L_\R$.
\end{proof}

\begin{lemma}[Gaussian part in the cone lies in one complex line]
\label{lem:gaussian-one-slope}
Let $\mu$ be an infinitely divisible probability measure on $\C^2$ supported on
the cone
\[
  \Gamma:=\{(u,v)\in\C^2:\ |u|=|v|\}.
\]
Assume that $\mu$ arises as a weak limit of sums of independent infinitesimal
vectors
\[
  V_{n,i}:=(a_{n,i}\xi_{n,i},\,b_{n,i}\xi_{n,i}),
\]
where each $\xi_{n,i}$ is centered, square-integrable, and satisfies
\[
  \gapgc(\xi_{n,i})>0.
\]
Then the Gaussian part of $\mu$ is supported on a single complex line
\[
  L_\tau:=\{(z,\tau z): z\in\C\}
  \qquad\text{for some } \tau\in\T.
\]
\end{lemma}

\begin{proof}
Let $G$ denote the Gaussian part of $\mu$, and let $E\subset \C^2\cong\R^4$ be
its support subspace.  Since $\mu$ is supported on $\Gamma$, the same is true of
$G$, so
\[
  E\subset \Gamma.
\]

The Gaussian covariance of the limit is obtained from the accumulation of the
row-wise covariance operators of the infinitesimal summands.  By
Lemma~\ref{lem:cov-range-rank-one}, every nonzero covariance contribution is the
real plane underlying a complex line
\[
  L_{n,i}:=\{(a_{n,i}z,b_{n,i}z): z\in\C\},
\]
because $\gapgc(\xi_{n,i})>0$ rules out support on a real line.

Hence the Gaussian support subspace $E$ is generated, as a real vector space, by
such planes $L_{n,i,\R}$.  Since $E\subset \Gamma$, every contributing plane is
itself contained in $\Gamma$.  Therefore, whenever $(a_{n,i},b_{n,i})\ne(0,0)$
and the corresponding covariance contribution is nonzero, one must have
\[
  |a_{n,i}|=|b_{n,i}|,
\]
so the associated complex line is of the form
\[
  L_{\tau_{n,i}}:=\{(z,\tau_{n,i}z): z\in\C\}
  \qquad\text{with } \tau_{n,i}\in\T.
\]

Assume for contradiction that two distinct slopes $\tau_1\ne\tau_2$ occur among
the planes contributing to $E$.  Then $E$ contains both complex lines
$L_{\tau_1}$ and $L_{\tau_2}$.  Since $E$ is a real vector subspace, it is
closed under addition, so
\[
  (L_{\tau_1}\cup L_{\tau_2}) + (L_{\tau_1}\cup L_{\tau_2}) \subset E\subset \Gamma.
\]
Applying Lemma~\ref{lem:two-slopes-real-degenerate} with
\[
  S:=L_{\tau_1}\cup L_{\tau_2}
\]
shows that the scalar set on each slope must be contained in a real line
through the origin.  This is impossible because each contributing set is the
full complex plane of scalars.  Hence only one slope can occur.

Therefore all covariance contributions lie in one complex line $L_\tau$, and so
their sum, namely the Gaussian covariance of $\mu$, has range contained in the
underlying real plane $L_{\tau,\R}$.  Thus $G$ is supported on $L_\tau$.
\end{proof}

\begin{lemma}[Support of the L\'evy measure inherits cone additivity]
\label{lem:levy-cone-additivity}
Let $\mu$ be an infinitely divisible probability measure on $\C^2$ with L\'evy
measure $\nu$, and assume that
\[
  \supp(\mu)\subset \Gamma:=\{(u,v)\in\C^2:\ |u|=|v|\}.
\]
Then:
\begin{enumerate}[label=\rm(\arabic*)]
\item $\supp(\nu)\subset \Gamma\cup\{0\}$;
\item if $x,y\in \supp(\nu)\setminus\{0\}$, then
\[
  x+y\in \Gamma.
\]
\end{enumerate}
\end{lemma}

\begin{proof}
Fix $x\in \supp(\nu)\setminus\{0\}$.  Choose an open neighborhood $U$ of $x$
whose closure is bounded away from the origin.  Then $0<\nu(U)<\infty$.  Write
\[
  \nu = \nu|_U + \nu|_{U^c}
\]
and let $\mu_U$ and $\mu_{U^c}$ be the corresponding infinitely divisible laws,
so that
\[
  \mu = \mu_U * \mu_{U^c}.
\]
Since $\nu(U)<\infty$, the law $\mu_U$ is a compound Poisson law.  In
particular, with positive probability it has exactly one jump in $U$ and no
other jumps, so every point of $U$ belongs to $\supp(\mu_U)$.  As convolution
supports add, this implies
\[
  U \subset \supp(\mu_U)\subset \supp(\mu)-\supp(\mu_{U^c}).
\]
Taking closures and using $\supp(\mu)\subset \Gamma$ yields $x\in \Gamma$.
This proves \textup{(1)}.

Now let $x,y\in \supp(\nu)\setminus\{0\}$.  Choose disjoint open neighborhoods
$U,V$ of $x,y$ whose closures are bounded away from $0$.  Then
$0<\nu(U),\nu(V)<\infty$.  Let $\mu_{U,V}$ denote the compound Poisson law with
L\'evy measure $\nu|_{U\cup V}$.  With positive probability, a sample from
$\mu_{U,V}$ has exactly one jump in $U$, exactly one jump in $V$, and no other
jumps in $U\cup V$.  Hence
\[
  U+V \subset \supp(\mu_{U,V}).
\]
As before, since $\mu=\mu_{U,V}*\mu_{(U\cup V)^c}$ and
$\supp(\mu)\subset \Gamma$, it follows that every point of $U+V$ belongs to
$\Gamma$.  Letting the neighborhoods shrink to $x$ and $y$ gives
\[
  x+y\in \Gamma.
\]
\end{proof}

\begin{corollary}[Two slopes in the L\'evy measure force real-line degeneracy]
\label{cor:levy-two-slopes}
Let $\mu$ be an infinitely divisible probability measure on $\C^2$ supported on
\[
  \Gamma:=\{(u,v)\in\C^2:\ |u|=|v|\},
\]
and let $\nu$ be its L\'evy measure.  Suppose $\nu$ charges two distinct
unimodular slope lines
\[
  L_{\tau_1}=\{(z,\tau_1 z): z\in\C\},
  \qquad
  L_{\tau_2}=\{(z,\tau_2 z): z\in\C\},
  \qquad
  \tau_1\ne \tau_2.
\]
Define
\[
  K_j:=\{z\in\C:\ (z,\tau_j z)\in \supp(\nu)\}
  \qquad (j=1,2).
\]
Then every nonzero point of $K_1$ forces $K_2$ into a real line through the
origin, and conversely.
\end{corollary}

\begin{proof}
Apply Lemma~\ref{lem:levy-cone-additivity} with
\[
  S:=\supp(\nu)\setminus\{0\}.
\]
Then $S\subset \Gamma$ and $S+S\subset \Gamma$.  Lemma~\ref{lem:two-slopes-real-degenerate}
therefore applies.
\end{proof}

\begin{remark}
Corollary~\ref{cor:levy-two-slopes} is the jump-part analogue of
Lemma~\ref{lem:two-slopes-real-degenerate}.  It shows that the only possible
way for the L\'evy measure to charge two distinct unimodular slopes is through a
real-line degeneracy of the scalar jump sets on those slopes.
\end{remark}

\begin{lemma}[Cone-supported infinitely divisible laws live on additive subgroups]
\label{lem:id-cone-additive-subgroup}
Let $\mu$ be an infinitely divisible probability measure on
\[
  \Gamma:=\{(u,v)\in\C^2:\ |u|=|v|\}.
\]
Then there exists an additive subgroup $G\subset \C^2$ such that
\[
  \supp(\mu)\subset \overline{G}\subset \Gamma.
\]
\end{lemma}

\begin{proof}
Since $\mu$ is infinitely divisible, there exists a probability measure $\nu$
such that
\[
  \mu=\nu*\nu.
\]
Let
\[
  A:=\supp(\nu).
\]
By the support formula for convolutions on $\R^4\cong\C^2$,
\[
  \supp(\mu)=\supp(\nu*\nu)=\overline{A+A}.
\]
Since $\supp(\mu)\subset \Gamma$ and $\Gamma$ is closed, it follows that
\[
  A+A\subset \Gamma.
\]

Define the real quadratic form
\[
  Q(u,v):=|u|^2-|v|^2
\]
and the associated symmetric real bilinear form
\[
  B\bigl((u,v),(u',v')\bigr):=\Re(u\overline{u'}-v\overline{v'}).
\]
Then
\[
  Q(x+y)=Q(x)+Q(y)+2B(x,y).
\]
If $x\in A$, then $2x=x+x\in A+A\subset \Gamma$, so
\[
  0=Q(2x)=4Q(x),
\]
hence $x\in \Gamma$.  Thus
\[
  A\subset \Gamma.
\]
Now if $x,y\in A$, then $x+y\in A+A\subset \Gamma$, and therefore
\[
  0=Q(x+y)=Q(x)+Q(y)+2B(x,y)=2B(x,y),
\]
because $Q(x)=Q(y)=0$.  Hence
\[
  B(x,y)=0
  \qquad (x,y\in A).
\]

Let $G$ be the additive subgroup generated by $A$.  Thus every $g\in G$ has
the form
\[
  g=\sum_{j=1}^m n_j a_j,
  \qquad n_j\in\mathbb Z,\ a_j\in A.
\]
Using bilinearity and the pairwise orthogonality just proved,
\[
  Q(g)=\sum_{j=1}^m n_j^2 Q(a_j)+2\sum_{1\le i<j\le m} n_i n_j B(a_i,a_j)=0.
\]
So
\[
  G\subset \Gamma.
\]
Finally, since $A+A\subset G$,
\[
  \supp(\mu)=\overline{A+A}\subset \overline{G}\subset \Gamma.
\]
\end{proof}

\begin{lemma}[Additive subgroups of the cone]
\label{lem:additive-subgroups-cone}
Let $G\subset \Gamma$ be an additive subgroup, where
\[
  \Gamma:=\{(u,v)\in\C^2:\ |u|=|v|\}.
\]
Then there exists $\lambda\in\T$ such that
\[
  G\subset P_\lambda^+
  \qquad\text{or}\qquad
  G\subset P_\lambda^-,
\]
where
\[
  P_\lambda^+:=\{(z,\lambda z): z\in\C\},
  \qquad
  P_\lambda^-:=\{(z,\lambda \overline z): z\in\C\}.
\]
\end{lemma}

\begin{proof}
Let
\[
  H:=\operatorname{span}_{\R}(G)\subset \C^2\cong \R^4.
\]
Using the same forms
\[
  Q(u,v):=|u|^2-|v|^2,
  \qquad
  B\bigl((u,v),(u',v')\bigr):=\Re(u\overline{u'}-v\overline{v'}),
\]
we have, for every $x,y\in G$,
\[
  0=Q(x+y)-Q(x)-Q(y)=2B(x,y),
\]
because $x,y,x+y\in \Gamma$.  Thus
\[
  B(x,y)=0
  \qquad (x,y\in G).
\]
By bilinearity, $B$ vanishes on $H\times H$, so $H$ is a totally isotropic
real subspace for the form $B$.  Since $B$ has signature $(2,2)$ on $\R^4$,
one has
\[
  \dim_{\R} H\le 2.
\]

If $\dim_{\R}H\le 1$, then every nonzero vector of $H$ belongs to $\Gamma$, and
the line $H$ is contained in both some $P_\lambda^+$ and some $P_\lambda^-$.
So the conclusion is immediate.

Assume now that $\dim_{\R}H=2$.  Consider the projection
\[
  \pi_1:H\to\C,
  \qquad
  \pi_1(u,v):=u.
\]
If $(0,v)\in H$, then
\[
  |v|=|0|=0,
\]
so $(0,v)=0$.  Thus $\pi_1$ is injective.  Since both $H$ and $\C$ are
2-dimensional over $\R$, the map $\pi_1$ is a real-linear isomorphism.  Hence
$H$ is the graph of a real-linear map $T:\C\to\C$:
\[
  H=\{(z,Tz): z\in\C\}.
\]
Because $H\subset \Gamma$, one has
\[
  |Tz|=|z|
  \qquad (z\in\C),
\]
so $T$ is a real-linear isometry of $\C\cong \R^2$.  Therefore either
\[
  Tz=\lambda z
  \qquad\text{for all } z\in\C,
\]
or
\[
  Tz=\lambda \overline z
  \qquad\text{for all } z\in\C,
\]
for some $\lambda\in\T$.  Thus
\[
  H=P_\lambda^+
  \qquad\text{or}\qquad
  H=P_\lambda^-.
\]
Since $G\subset H$, the proof is complete.
\end{proof}

\begin{remark}
Lemmas~\ref{lem:id-cone-additive-subgroup} and
\ref{lem:additive-subgroups-cone} give an alternative structural description of
cone-supported infinitely divisible laws: their support is contained in the
closure of an additive subgroup lying either in a complex line
$P_\lambda^+$ or in a conjugate plane $P_\lambda^-$.  This is consistent with
the Gaussian obstruction $(Z,\overline Z)$ and explains why the rank-one
structure of the array is needed to rule out the conjugate-plane alternative.
\end{remark}

\begin{lemma}[Diffuse tails force a nontrivial Gaussian part]
\label{lem:tail-gaussian-nontrivial}
Let
\[
  V_{n,i}:=(a_i^{(n)}\xi_{n,i},\,b_i^{(n)}\xi_{n,i})\in\C^2\cong\R^4
\]
for $i>m_n$, where the coordinates $\xi_{n,i}$ are independent and satisfy
\[
  \E[\xi_{n,i}]=0,
  \qquad
  \E[|\xi_{n,i}|^2]=1,
  \qquad
  \E[|\xi_{n,i}|]\ge A>0.
\]
Define
\[
  \alpha_{n,i}:=\bigl(|a_i^{(n)}|^2+|b_i^{(n)}|^2\bigr)^{1/2}.
\]
Assume that
\[
  \max_{i>m_n}\alpha_{n,i}\longrightarrow 0,
  \qquad
  \sum_{i>m_n}\alpha_{n,i}^2\longrightarrow \sigma^2>0,
\]
and that the row sums
\[
  \sum_{i>m_n}V_{n,i}
\]
converge in distribution to an infinitely divisible law $\mu$ on $\C^2$.
Then the Gaussian part of $\mu$ is nontrivial.
\end{lemma}

\begin{proof}
Let
\[
  \beta_n:=\max_{i>m_n}\alpha_{n,i},
  \qquad
  r_n:=\beta_n^{1/2}.
\]
Then
\[
  r_n\downarrow 0,
  \qquad
  \frac{r_n}{\beta_n}\longrightarrow \infty.
\]
For each $i>m_n$ with $\alpha_{n,i}\ne 0$, set
\[
  K_{n,i}:=\frac{r_n}{\alpha_{n,i}}\ge \frac{r_n}{\beta_n}=\beta_n^{-1/2},
\]
so $K_{n,i}\to\infty$ uniformly in $i>m_n$.  Define the centered truncation
\[
  \widetilde\xi_{n,i}
  :=
  \xi_{n,i}\mathbf 1_{\{|\xi_{n,i}|\le K_{n,i}\}}
  -
  \E\!\left[\xi_{n,i}\mathbf 1_{\{|\xi_{n,i}|\le K_{n,i}\}}\right],
\]
and set
\[
  \widetilde V_{n,i}:=(a_i^{(n)}\widetilde\xi_{n,i},\,b_i^{(n)}\widetilde\xi_{n,i}).
\]
If $\alpha_{n,i}=0$, simply set $\widetilde\xi_{n,i}=0$ and
$\widetilde V_{n,i}=0$.

For every $K>0$ and every centered $\xi$ with $\E|\xi|^2=1$ and
$\E|\xi|\ge A$, one has
\[
  \E|\xi|
  \le
  \left(\E\!\left[|\xi|^2\mathbf 1_{\{|\xi|\le K\}}\right]\right)^{1/2}
  +
  \E\!\left[|\xi|\mathbf 1_{\{|\xi|>K\}}\right]
  \le
  \left(\E\!\left[|\xi|^2\mathbf 1_{\{|\xi|\le K\}}\right]\right)^{1/2}
  +
  \frac1K,
\]
and therefore
\[
  \E\!\left[|\xi|^2\mathbf 1_{\{|\xi|\le K\}}\right]
  \ge
  \left(A-\frac1K\right)^2.
\]
Also, since $\E[\xi]=0$,
\[
  \left|
    \E\!\left[\xi\mathbf 1_{\{|\xi|\le K\}}\right]
  \right|
  =
  \left|
    \E\!\left[\xi\mathbf 1_{\{|\xi|>K\}}\right]
  \right|
  \le
  \E\!\left[|\xi|\mathbf 1_{\{|\xi|>K\}}\right]
  \le
  \frac1K.
\]
Hence
\[
  \operatorname{Var}\!\left(
    \xi\mathbf 1_{\{|\xi|\le K\}}
    -
    \E\!\left[\xi\mathbf 1_{\{|\xi|\le K\}}\right]
  \right)
  \ge
  \left(A-\frac1K\right)^2-\frac1{K^2}.
\]
Choose $K_0$ large enough that
\[
  c_0:=\left(A-\frac1{K_0}\right)^2-\frac1{K_0^2}>0.
\]
Since $K_{n,i}\to\infty$ uniformly, for all sufficiently large $n$ one has
$K_{n,i}\ge K_0$ for every $i>m_n$.  Therefore
\[
  \E|\widetilde\xi_{n,i}|^2\ge c_0
  \qquad (i>m_n).
\]

Next,
\[
  |\widetilde V_{n,i}|
  \le
  |V_{n,i}|\mathbf 1_{\{|V_{n,i}|\le r_n\}}
  +
  \E\!\left[|V_{n,i}|\mathbf 1_{\{|V_{n,i}|\le r_n\}}\right]
  \le
  2r_n,
\]
so the centered truncation array is bounded and infinitesimal.  Moreover,
\[
  \E|\widetilde V_{n,i}|^2
  =
  \alpha_{n,i}^2\,\E|\widetilde\xi_{n,i}|^2
  \ge
  c_0\,\alpha_{n,i}^2
\]
for all large $n$.  Summing over $i>m_n$ gives
\[
  \sum_{i>m_n}\E|\widetilde V_{n,i}|^2
  \ge
  c_0\sum_{i>m_n}\alpha_{n,i}^2
  \longrightarrow
  c_0\sigma^2>0.
\]

By the standard shrinking-truncation theorem for infinitesimal triangular
arrays and infinitely divisible limits, the covariance matrices of these
centered truncations converge to the Gaussian covariance matrix of $\mu$; see
for instance Parthasarathy~\cite[Chapter~IV, \S5, pp.~82--96; \S6,
pp.~97--101; \S7, pp.~102--108]{parthasarathy}; these are the passages
where the accompanying infinitely divisible laws are constructed and the
Gaussian component is identified through the covariance of shrinking centered
truncations.  Since the traces stay bounded away from zero, that Gaussian
covariance matrix is nonzero.  Thus the Gaussian part of $\mu$ is nontrivial.
\end{proof}

\begin{lemma}[Degenerate cone case reduces to a line]
\label{lem:degenerate-cone-line-reduction}
Let $(R_X,R_Y)$ be the tail limit given by Lemma~\ref{lem:limit-law}, and
assume
\[
  H_X\equiv 0,
  \qquad
  H_Y\equiv 0.
\]
Then there exists $\tau\in\T$ such that the law of
\[
  R_Y-\tau R_X
\]
is supported on an affine line.
\end{lemma}

\begin{proof}
Let
\[
  \mu:=\text{law}(R_X,R_Y).
\]
Since $H_X\equiv H_Y\equiv 0$, Lemma~\ref{lem:limit-law} gives
\[
  |R_X|=|R_Y|
  \qquad\text{almost surely,}
\]
so
\[
  \supp(\mu)\subset \Gamma:=\{(u,v)\in\C^2:\ |u|=|v|\}.
\]

By Lemmas~\ref{lem:id-cone-additive-subgroup} and
\ref{lem:additive-subgroups-cone}, there exists $\lambda\in\T$ such that
\[
  \supp(\mu)\subset P_\lambda^+
  \qquad\text{or}\qquad
  \supp(\mu)\subset P_\lambda^-,
\]
where
\[
  P_\lambda^+:=\{(z,\lambda z): z\in\C\},
  \qquad
  P_\lambda^-:=\{(z,\lambda \overline z): z\in\C\}.
\]

If $\supp(\mu)\subset P_\lambda^+$, then
\[
  R_Y=\lambda R_X
  \qquad\text{almost surely,}
\]
so choosing $\tau:=\lambda$ gives
\[
  R_Y-\tau R_X=0
  \qquad\text{almost surely.}
\]
This is already stronger than the conclusion.

It remains to treat the case
\[
  \supp(\mu)\subset P_\lambda^-.
\]
Then for every support point of $\mu$ there exists $z\in\C$ such that
\[
  (R_X,R_Y)=(z,\lambda\overline z).
\]
Choosing $\tau:=\lambda$, we obtain
\[
  R_Y-\tau R_X
  =
  \lambda(\overline z-z)
  \in
  i\lambda\,\R.
\]
Thus the law of $R_Y-\tau R_X$ is supported on the real line
\[
  i\lambda\,\R\subset\C.
\]
\end{proof}

\begin{remark}
Lemma~\ref{lem:tail-gaussian-nontrivial} is exactly where the uniform lower
$L^1$ bound enters the degenerate-head case.  It gives a positive lower bound
for the variances of shrinking centered truncations without any appeal to
uniform $L^2$-integrability of the coordinate family.
\end{remark}

\begin{remark}
The jump-part lemmas above remain useful as complementary structural
information.  In particular, Corollary~\ref{cor:levy-two-slopes} shows that any
attempt to prove the theorem using only the L\'evy measure would still be
forced into a real-line-degenerate two-slope analysis.
\end{remark}

\begin{remark}[Historical route via quadratic mass]
Before Lemma~\ref{lem:tail-gaussian-nontrivial} was isolated, the degenerate
case was naturally led by the observation
\[
  \sum_{i>m_n}\bigl(|a_i^{(n)}|^2+|b_i^{(n)}|^2\bigr)\longrightarrow 2.
\]
Lemma~\ref{lem:tail-gaussian-nontrivial} is the formal version of that
heuristic: the shrinking-truncation method does produce a nonzero Gaussian part
once the uniform lower $L^1$ bound is fed into the truncation estimates.
\end{remark}

\section{Proof of Theorem~\ref{thm:main}}

\begin{proof}[Proof of Theorem~\ref{thm:main}]
Assume for contradiction that Theorem~\ref{thm:main} fails.  Let the bad
sequence be given by Proposition~\ref{prop:counterexample}.  Apply
Proposition~\ref{prop:compactness-reduction}.  Thus, after passing to a
subsequence, there exist head and tail limits $(H_X,H_Y,R_X,R_Y)$ with
$(H_X,H_Y)$ independent of $(R_X,R_Y)$ and
\[
  |H_X+R_X|=|H_Y+R_Y|
  \qquad\text{almost surely.}
\]

\medskip
\noindent
\textbf{Case 1: $H_X$ is not almost surely constant.}
By Proposition~\ref{prop:compactness-reduction}, there exists $\tau\in\T$ such
that
\[
  H_Y=\tau H_X
  \qquad\text{almost surely.}
\]
Since $\supp(H_X)$ contains at least two distinct points,
Lemma~\ref{lem:nontrivial-head-geometry} shows that the law of
\[
  W:=R_Y-\tau R_X
\]
is supported on an affine line, hence on a generalized circle.

Now
\[
  W_n:=Y_n^{\mathrm{tail}}-\tau X_n^{\mathrm{tail}}
  =
  \sum_{i>m_n}(b_i^{(n)}-\tau a_i^{(n)})\,\xi_{n,i}
\]
converges in distribution to $W$.  The coefficients are diffuse because
Lemma~\ref{lem:profile} gives
\[
  \max_{i>m_n}|b_i^{(n)}-\tau a_i^{(n)}|
  \le
  2\max_{i>m_n}\max\bigl(|a_i^{(n)}|,|b_i^{(n)}|\bigr)\longrightarrow 0.
\]
Therefore Proposition~\ref{prop:tail-rigidity} implies
\[
  \|b^{(n)}_{\mathrm{tail}}-\tau a^{(n)}_{\mathrm{tail}}\|_2\longrightarrow 0.
\]
By Proposition~\ref{prop:head},
\[
  \|b^{(n)}_{\mathrm{head}}-\tau a^{(n)}_{\mathrm{head}}\|_2\longrightarrow 0.
\]
Hence
\[
  \|b^{(n)}-\tau a^{(n)}\|_2\longrightarrow 0.
\]
This contradicts the orthogonal normalization, since
\[
  \|b^{(n)}-\tau a^{(n)}\|_2^2
  =
  \|a^{(n)}\|_2^2+\|b^{(n)}\|_2^2
  -2\Re\bigl(\tau\langle a^{(n)},b^{(n)}\rangle\bigr)
  = 2.
\]

\medskip
\noindent
\textbf{Case 2: $H_X$ is almost surely constant.}
Since $H_X$ is centered as a weak limit of centered random variables, one has
\[
  H_X\equiv 0,
  \qquad
  H_Y\equiv 0.
\]
By Lemma~\ref{lem:degenerate-cone-line-reduction}, there exists $\tau_0\in\T$
such that the law of
\[
  W:=R_Y-\tau_0 R_X
\]
is supported on an affine line, hence on a generalized circle.  Therefore
Proposition~\ref{prop:tail-rigidity} applied to
\[
  W_n:=Y_n^{\mathrm{tail}}-\tau_0 X_n^{\mathrm{tail}}
\]
shows that
\[
  \|b^{(n)}_{\mathrm{tail}}-\tau_0 a^{(n)}_{\mathrm{tail}}\|_2\longrightarrow 0.
\]
Because $H_X\equiv H_Y\equiv 0$, the head limits vanish, so
\[
  \|a^{(n)}_{\mathrm{head}}\|_2\longrightarrow 0,
  \qquad
  \|b^{(n)}_{\mathrm{head}}\|_2\longrightarrow 0.
\]
Thus
\[
  \|b^{(n)}-\tau_0 a^{(n)}\|_2\longrightarrow 0,
\]
again contradicting the orthogonal normalization.
\end{proof}

\begin{remark}[Alternative Gaussian route in the degenerate case]
The proof of Lemma~\ref{lem:degenerate-cone-line-reduction} does not use the
Gaussian part.  Nevertheless, the same conclusion also follows from
Lemma~\ref{lem:tail-gaussian-nontrivial}.  Indeed, if
\[
  |R_X|=|R_Y|
\]
and one writes the tail law as a convolution of its Gaussian part and its
independent remainder, then the Gaussian part is nontrivial and, by
Lemma~\ref{lem:gaussian-one-slope}, lies in a single complex line
\[
  \{(z,\tau z): z\in\C\}.
\]
Using the exact modulus identity and Lemma~\ref{lem:nontrivial-head-geometry}
with the Gaussian first coordinate as the new head again forces
$R_Y-\tau R_X$ onto an affine line.  Thus the degenerate case admits both a
pure support proof and an independent Gaussian-based proof.
\end{remark}

\bibliographystyle{plain}
\begin{thebibliography}{10}

\bibitem{complex_spr}
P.~Abdalla Teixeira, J.~de~Dios Pont, J.~P.~Ramos, and M.~Taylor,
\newblock A complex stable phase retrieval theorem,
\newblock Preprint, 2024.

\bibitem{complex_spr_app}
P.~Abdalla Teixeira, J.~de~Dios Pont, J.~P.~Ramos, and M.~Taylor,
\newblock A complex stable phase retrieval theorem, appendix draft,
\newblock Preprint, 2024.

\bibitem{compactness_real}
P.~Abdalla Teixeira, J.~de~Dios Pont, J.~P.~Ramos, and M.~Taylor,
\newblock A compactness proof of stable phase retrieval,
\newblock Preprint, 2024.

\bibitem{FOTP}
D.~Freeman, D.~Okpoti, S.~A.~Teixeira, and D.~Phong,
\newblock Phase retrieval in infinite-dimensional Hilbert spaces,
\newblock Preprint, 2023.

\bibitem{calderbank2022stable}
A.~R.~Calderbank, I.~Daubechies, B.~Freeman, and S.~Freeman,
\newblock Stable phase retrieval and perturbations of frames,
\newblock {\em SIAM J.\ Appl.\ Algebra Geom.}, 2022.

\bibitem{gnedenko-kolmogorov}
B.~V.~Gnedenko and A.~N.~Kolmogorov,
\newblock {\em Limit Distributions for Sums of Independent Random Variables},
\newblock Addison-Wesley, 1954.

\bibitem{parthasarathy}
K.~R.~Parthasarathy,
\newblock {\em Probability Measures on Metric Spaces},
\newblock Academic Press, 1967; AMS Chelsea reprint, 2005.

\end{thebibliography}

\end{document}
